\section{Haar measurement simulation and the geometry of contact}
\label{geo:qutrit}

The same projective moments have an operational use beyond ensemble
certification. A \emph{parent measurement} is a fixed measurement from
which a family of measurements can be obtained by classical
postprocessing. In dimension \(d\), the unitary-invariant pure-state
parent is
\[
 G(dz)=dP_z\,d\sigma(z).
\]
For \(m\) labelled outcomes its postprocessing body is
\[
 \mathcal H_{d,m}=
 \left\{\left(d\int p_i(z)P_z\,d\sigma(z)\right)_{i=1}^m:
      p_i\ge0,\quad\sum_i p_i=1\ {\rm a.e.}\right\}.
\]
It is compact and convex by weak-star compactness of bounded response
functions. Its support function,
\[
 h(W)=d\int\max_i\operatorname{tr}(W_iP_z)\,d\sigma(z),
 \label{geo:haar-support}
\]
turns simulation into a comparison between a classical maximum and
quantum measurement optimization. This convex-body description belongs
to the established theory of steering and measurement compatibility
\cite{geo_nguyen2018}.

Let
\[
 \Phi_\eta(A)=\eta A+\frac{1-\eta}{d}\operatorname{tr}(A)I,\qquad
 \eta_d=\frac{H_d-1}{d-1},\qquad c_d=\frac{1-\eta_d}{d}.
\]
The notation \(d\) in this section denotes Hilbert-space dimension.
The number \(\eta_d\) is fixed by the exact maximum-coordinate cells
\(R_a=\{z:|z_a|^2\ge|z_b|^2\text{ for every }b\}\):
\begin{equation}
 d\int_{R_a}P_z\,d\sigma(z)=\eta_d P_a+c_dI.
 \label{geo:haar-cell}
\end{equation}
Indeed, \((|z_1|^2,\ldots,|z_d|^2)\) is uniform on the simplex and
\(\mathbb E\max_a|z_a|^2=H_d/d\).
Thus the same expected-maximum computation governing the entropy
threshold also determines the noise level at which every basis
measurement is a postprocessing of one Haar parent.

\fig{fig10_measurement_cells}{The Haar parent at qutrit visibility $5/12$.
(a) Squared coordinate moduli of a Haar qutrit are uniform on the simplex.
Assigning each point to its largest coordinate partitions it into three
cells; phase integration removes off-diagonal matrix entries. (b) The three
postprocessed effects have eigenvalues $11/18$ in their selected direction
and $7/36$ in the other two directions. Segment heights are the three
diagonal eigenvalues, colored by coordinate direction; each complete column
sums to one. Hence each equals
$(5/12)P_a+(7/36)I$. The neighborhood theorem extends this exact basis-cell
construction to the specified classes of nearby measurements; the figure
itself depicts the basis case.}{fig:measurement}

Write \(\mathsf Q_{d,m}\) for the compact set of labelled \(m\)-outcome
POVMs, and \(\mathsf P_{d,m}\) for the convex hull of its commuting
members. This is the standard projective-simulable set
\cite{geo_oszmaniec2017}.
For qutrits, let \(\mathsf R_m\) consist of all POVMs with a common
nontrivial reducing subspace, and set
\(\mathsf K_m=\operatorname{conv}\mathsf R_m\). A qutrit reducing
decomposition is a two-dimensional block plus an orthogonal flag.
The decomposition can vary between convex-combination terms.

\begin{theorem}[Uniform operational neighborhoods]
\label{geo:operational}
For every fixed \(d\ge2\) and finite \(m\), there exists \(r_{d,m}>0\)
such that
\[
 E\in\mathsf Q_{d,m},\quad
 \operatorname{dist}(E,\mathsf P_{d,m})<r_{d,m}
 \quad\Longrightarrow\quad
 \Phi_{\eta_d}(E)\in\mathcal H_{d,m}.
 \tag{i}
\]
For \(d=3\), there exists \(r_m>0\) with the stronger conclusion
\[
 E\in\mathsf Q_{3,m},\quad
 \operatorname{dist}(E,\mathsf K_m)<r_m
 \quad\Longrightarrow\quad
 \Phi_{5/12}(E)\in\mathcal H_{3,m}.
 \tag{ii}
\]
Any fixed norm on tuples can define distance. Zero labelled effects
are allowed, and nearby POVMs may be irreducible and have arbitrary
effect ranks. The radii are qualitative; no value uniform in
unbounded \(m\) is asserted.
\end{theorem}

In particular, every reducible qutrit POVM and every convex mixture of
such measurements is simulated by the same Haar parent at \(5/12\).
Every qutrit POVM with at most three outcomes lies in this class.
To check the last statement, two effects with overlapping support
admit opposite small Hermitian perturbations supported in their
intersection, contradicting extremality. In dimension three, an
extreme ternary POVM therefore has either three rank-one effects,
hence is projective, or rank pattern \((2,1,1)\), with smaller cases
allowed. In the latter case choose a ray \(f\) orthogonal to the
supports of both rank-one effects. Those effects annihilate \(f\),
and the remaining effect, being their complement in \(I\), satisfies
\(E_1f=f\). This is a common reducing flag. The finite-dimensional
convex decomposition into extreme POVMs finishes the assertion.

\begin{remark}
The theorem is an open-neighborhood result around a specified compact
class. It does not assert that all irreducible qutrit POVMs are
simulable at \(5/12\). The contact theorem below is fully anisotropic
on the flagged stratum: its restriction is geometric support, not a
restriction on the spectra of the contact matrices.
\end{remark}

\subsection{Contact inequalities and martingale transport}
\label{geo:contact}

Quantum discrimination gives the semidefinite duality
\[
 q(W):=\max_{E\in\mathsf Q_{d,m}}
              \sum_i\operatorname{tr}(\Phi_{\eta_d}(W_i)E_i)
       =\min_{Y\succeq\Phi_{\eta_d}(W_i)}\operatorname{tr}Y.
\]
For a dual optimum put \(C_i=Y-\Phi_{\eta_d}(W_i)\succeq0\).
Complementary slackness says that an active rank-one effect lies in
\(\ker C_i\); its weighted rays are balanced. Define
\[
 F_d(C)=\int\min_i
       \{z^*C_i z-c_d\operatorname{tr}C_i\}\,d\sigma(z).
\]
Since
\(\Phi_{\eta_d}^{-1}(C)=
\eta_d^{-1}(C-c_d\operatorname{tr}(C)I)\), one has the exact identity
\begin{equation}
 h(W)=\operatorname{tr}Y-\frac d{\eta_d}F_d(C).
 \label{geo:contact-identity}
\end{equation}
Thus a nonpositive contact integral controls the support comparison.
The general contact formulation appears in
\cite[Corollary 2, Eq.~(35)]{geo_yzhang2026}; the identity above also
provides its direct derivation in the present normalization.

There is a second useful interpretation. For a rank-one POVM
\(E_i=t_iP_i\), let \(w_i=t_i/d\) and
\(\rho_i=\eta P_i+(1-\eta)I/d\).
A Haar postprocessing of \(\Phi_\eta(E)\) is exactly a joint law of
the label \(i\) and a Haar pure projector \(P\) satisfying
\[
 \mathbb P(i)=w_i,\qquad \mathbb E[P\mid i]=\rho_i.
\]
Consequently, by the classical martingale coupling theorem
\cite{geo_strassen1965}, simulation is equivalent to
\[
 \sum_iw_i\delta_{\rho_i}\preceq_{\rm cx}\sigma
\]
in the real affine space of trace-one Hermitian matrices.
This is ordinary classical convex order on a matrix-valued state
space. The contact functions are specific convex-envelope tests for
that order; testing only covariance or quadratic polynomials is weaker.

\subsection{The anisotropic qutrit flag theorem}
\label{geo:flag}

\begin{theorem}[Arbitrary contacts on a flagged qutrit frame]
\label{geo:flag-theorem}
Let \(\mathbb C^3=H\oplus\mathbb Cf\), \(\dim H=2\), and let
\(P_i\) be finitely many rays in \(H\) with
\(\sum_i t_iP_i=I_H\), \(t_i\ge0\). Include \(P_f\).
For every collection of positive matrices \(C_i,C_f\) with
\(C_iP_i=0\) and \(C_fP_f=0\), one has
\[
 F_3(C)\le0,\qquad \eta_3=\frac5{12},\quad c_3=\frac7{36}.
\]
\end{theorem}

\begin{proof}
Average the contacts under the flag phase
\(I_H\oplus e^{i\chi}\). Concavity of \(F_3\) and Haar invariance
show that this can only increase \(F_3\), so it suffices to consider
\[
 C_i=a_i(I_H-P_i)+b_iP_f,\quad a_i,b_i\ge0,\qquad
 C_f\succeq0\text{ on }H.
\]
In Bloch notation \(P_i=(I_H+n_i\cdot\sigma)/2\), balance means
\(0\in\operatorname{conv}\{n_i\}\). Carath\'eodory reduces to at most
four rays: deleting contacts increases the pointwise minimum.
One or two rays give a basis case by \eqref{geo:haar-cell}.

For three selected rays \(i,j,k\), put
\[
 w=a_in_i+a_jn_j-2a_kn_k,\quad
 v=a_in_i-a_jn_j,\quad A=a_i+a_j+2a_k.
\]
Choose a qubit basis with Bloch axis \(n\). On its first
maximum-coordinate cell assign the better of contacts \(i,j\), on
the second assign \(k\), and on the flag cell assign \(f\).
The exact diagonal moments cancel the flag coefficients. Averaging
the absolute value of the pair difference over its relative phase
gives
\begin{equation}
 F_3(C)\le\frac5{144}
 \left(A-n\cdot w-t_*|v\times n|\right),
 \quad
 t_*=\frac{36K_0}{5},\quad
 K_0=\frac4\pi\int_{R_1}\sqrt{p_1p_2}\,d\sigma.
 \label{geo:pair-cell}
\end{equation}
The phase inequality is
\(\mathbb E_\phi|b+a\cos\phi|\ge2|a|/\pi\).
An elementary one-dimensional integration gives
\[
 \int_{R_1}\sqrt{p_1p_2}\,d\sigma
 =\frac23\int_{1/2}^1\sqrt{u(1-u)}
                       \bigl(1-(1+u)^{-3}\bigr)\,du
 \ge\frac{19\pi}{648}+\frac{37}{3888}.
\]
Using \(\pi<22/7\),
\begin{equation}
 t_*\ge\frac{2767}{2970}>\frac9{10},
 \qquad t_*-\frac9{10}\ge\frac{47}{1485}.
 \label{geo:strict-reserve}
\end{equation}
The needed geometric statement is that among any at most four
balanced Bloch rays, for arbitrary \(a_i\ge0\), some triple satisfies
\begin{equation}
 |w|^2+\frac{81}{100}|v|^2+
             \frac95|w\times v|\ge A^2.
 \label{geo:four-selection}
\end{equation}
Its proof is given below, including the nearest-facet estimates.
Finally,
\[
 \max_{|n|=1}\{n\cdot w+t|v\times n|\}
 =\sqrt{|w|^2+t^2|v|^2+2t|w\times v|}.
\]
Equations \eqref{geo:four-selection} and \eqref{geo:strict-reserve}
make the right side of \eqref{geo:pair-cell} nonpositive.
\end{proof}

% The geometric proof may instead be moved to a technical appendix.

For a reducible POVM, pinching any separating witness to the same
two-plus-one block leaves its score unchanged and decreases Haar
support. The corresponding quantum SDP splits into a qubit measurement
and assignment of the flag. An extreme qubit optimizer has at most
four rank-one effects, with a balanced set of Bloch rays.
Complementary slackness gives exactly the contacts in
Theorem~\ref{geo:flag-theorem}; \eqref{geo:contact-identity} rules out
separation. Convexity then proves
\(\Phi_{5/12}(\mathsf K_m)\subseteq\mathcal H_{3,m}\).

\subsection{Why local contact control implies an operational neighborhood}
\label{geo:op-proof}

The main distinction is between the rays of an external POVM and
the rays of an optimizing dual contact frame. A compactness argument
must bridge them. The following estimate supplies the required control.

\begin{lemma}[A negative cusp for pinched clusters]
\label{geo:cusp}
Suppose a balanced \(m\)-ray frame is partitioned into \(d\) nonempty
groups, with actual maximum within-group projector distance
\(\delta\le1/(2d)\). Suppose
\[
 a(I-P_i)\preceq C_i\preceq b(I-P_i),\qquad 0<a\le b.
\]
For a constant \(\kappa_d>0\) depending only on dimension,
\begin{equation}
 F_d(C)\le9\eta_db\delta^2-\frac{\kappa_da}{m}\delta.
 \label{geo:negative-cusp}
\end{equation}
One may take \(\kappa_3=325/536313\).
\end{lemma}

\begin{proof}
For group \(I_a\), form its weighted mean
\(A_a=T_a^{-1}\sum_{i\in I_a}t_iP_i\). Its largest eigenvalue is
at least \(1-\delta^2\), and its remaining spectral mass at most
\(\delta^2\). If \(Q_a=u_au_a^*\) is its principal projector,
\(\|P_i-Q_a\|\le\sqrt2\delta\).
Balance gives
\(\sum_aT_a\lambda_aQ_a=I-S\), where \(S\succeq0\) and
\(\operatorname{tr}S\le d\delta^2\). No group can have zero total
weight, since \(I-S\) is invertible.
The polar decomposition of the matrix with columns
\(\sqrt{T_a\lambda_a}u_a\) gives an orthonormal basis \(e_a\)
with \(\|P_i-P_{e_a}\|<3\delta\) in group \(a\).
Consequently \(e_a^*C_ie_a\le9b\delta^2\).

On basis cell \(R_a\), bound the minimum by the uniform average of
its group's contacts. Equation \eqref{geo:haar-cell} bounds the
total integrated average by \(9\eta_db\delta^2\).
The functional
\[
 D\longmapsto\int_{R_a}
       |z^*Dz-c_d\operatorname{tr}D|\,d\sigma
\]
is a norm on Hermitian matrices: if its polynomial vanishes on the
open cell, then \(D=c_d\operatorname{tr}(D)I\), and
\(1-dc_d=\eta_d>0\) forces \(D=0\).
It is therefore at least \(\kappa_d\|D\|\).
In a group attaining diameter \(\delta\), choose \(i,j\) attaining it.
The average minus the minimum is at least \(|f_i-f_j|/|I_a|\).
Moreover, a unit vector in \(\operatorname{ran}P_j\) gives
\(\|C_i-C_j\|\ge a\delta\). This yields the negative linear term.
If \(\delta=0\), the balanced frame is a basis refinement, and
\eqref{geo:haar-cell} gives \(F_d\le0\) directly.
\end{proof}

\paragraph{Normal cones of the Haar body.}
Suppose \(W\) supports \(\mathcal H_{d,m}\) at the noisy image of a
commuting POVM \(Q_i=\sum_a p_i(a)P_a\).
Its known response equals \(p_i(a)\) on \(R_a\).
Every positive branch must therefore attain the pointwise maximum
throughout that cell. Phase variation on interfaces \(p_a=p_b\)
and variation of coordinate moduli show that selected branches have
the canonical form
\begin{equation}
 W_{i(a)}=M+\lambda P_a,\qquad \lambda\ge0,
 \label{geo:normal}
\end{equation}
and every other branch lies below the envelope
\[
 z^*W_jz\le z^*Mz+\lambda\max_a|z_a|^2.
 \label{geo:envelope}
\]
Indeed, the difference of two selected branches must be
\(\lambda_{ab}(P_a-P_b)\); the cocycle relation makes all
\(\lambda_{ab}\) equal, with sign fixed by the maximizing cells.

An unused branch is controlled quantitatively at the limiting
normal. If \(z^*Az\le\max_a|z_a|^2\) for every unit vector, then
\(A\preceq I\), and averaging over flat-amplitude phase states gives
\(\operatorname{tr}A\le1\). Thus
\[
 (\eta_d+c_d)I-\Phi_{\eta_d}(A)
 =\eta_d(I-A)+c_d(1-\operatorname{tr}A)I\succeq0.
\]
This slack is singular only if \(A=P_a\).
In fact singularity forces \(\operatorname{tr}A=1\) and \(Av=v\);
the envelope makes \(v\) a basis vector \(e_a\).
The nonnegative function \(z^*(P_a-A)z\) on \(R_a\) has zero
integral by \eqref{geo:haar-cell}. It vanishes on an open set,
so \(A=P_a\).

\paragraph{Support reduction for qutrits.}
Every supporting witness at the noisy image of a point of
\(\mathsf K_m\) also supports at a commuting POVM.
First take a reducible point and pinch \(W\) to its block,
obtaining \(\overline W\). The unchanged score, the flagged theorem,
and convexity of Haar support give
\[
 h(W)=\operatorname{score}_{\overline W}(E)
 \le q(\overline W)\le h(\overline W)\le h(W).
\]
All inequalities are equalities. Choose an extreme qubit optimizer
for \(\overline W\). If it has one or two rays it is commuting.
If an active contact coefficient \(a_i=0\), that contact annihilates
the entire qubit block and a single block effect is also optimal.
In the remaining three- or four-ray case, all \(a_i>0\) and the rays
are distinct. The strict reserve \eqref{geo:strict-reserve} makes
the pair-cell bound strictly negative, since \(v\ne0\) and
\(\sqrt{|w|^2+t^2|v|^2+2t|w\times v|}\) is strictly increasing in \(t>0\).
Equation \eqref{geo:contact-identity} would give \(h>q\), a contradiction.
A commuting optimizer therefore exists; as it is fixed by the
pinching, it has the same score for the original \(W\).
For a convex combination of reducible POVMs, a vanishing support
deficit forces every positive component's deficit to vanish.
This proves the reduction for \(\mathsf K_m\).
For \(\mathsf P_{d,m}\), the same reduction follows immediately
by taking a commuting component.

\begin{proof}[Proof of Theorem~\ref{geo:operational}]
Both center sets are compact and closed under merging labels and
deleting zero effects. We argue by induction on \(m\), beginning
with the single effect \(I\).
If nonsimulable \(E^n\) approached a center point \(E^0\), choose
separating witnesses \(W^n\), subtract their common matrix average,
and normalize their tuple norm to one. A subsequence converges to
a nongauge support \(W^0\) at \(E^0\).
The support reduction supplies a commuting support point, so
\eqref{geo:normal} and the envelope apply.

If \(\lambda>0\), the unique optimal quantum dual at the limit is
\[
 Y^0=\Phi_{\eta_d}(M)+\lambda(\eta_d+c_d)I.
\]
Feasibility follows from the envelope lemma. The selected basis
labels force every basis diagonal of an optimal dual to attain its
minimal value; positivity then kills all off-diagonal entries,
proving uniqueness.
Nearby optimal duals \(Y^n\) are bounded and converge to \(Y^0\).
Every limiting slack is either positive definite or
\(\eta_d\lambda(I-P_a)\).
Complementary slackness for optimal quantum POVMs makes the active
rays fall into \(d\) shrinking clusters about the basis, with
uniform positive pinching on their kernel complements. Every basis
group is present, since the active effects sum to \(I\).
Discarding inactive labels decreases Haar support without changing
the quantum value. Strict separation and
\eqref{geo:contact-identity} then force \(F_d(C^n)>0\).
This contradicts Lemma~\ref{geo:cusp} as the actual cluster diameter
tends to zero.

If \(\lambda=0\), all selected branches equal \(M\), while
\(W_j^0\preceq M\). Nongauge normalization provides a strict branch
\(D=M-W_j^0\succeq0\), \(D\ne0\), for which
\(\Phi_{\eta_d}(D)\succ0\). Equality of the limiting score forces
\(E_j^0=0\). For large \(n\), merge that label into a selected branch:
the quantum score does not decrease, while deleting the witness
branch does not increase Haar support. Strict separation survives
with \(m-1\) labels. The limit is a center point with its zero label
deleted, contradicting induction.

Thus no nonsimulable sequence can approach any point of the compact
center set. Compactness gives a uniform positive distance for each
fixed \(d,m\) (or fixed \(m\) in the qutrit case).
\end{proof}

\begin{corollary}[Exact boundary within the center classes]
\label{geo:boundary}
After deleting zero effects, a POVM in \(\mathsf P_{d,m}\), or in
\(\mathsf K_m\) for \(d=3\), maps to the boundary of the Haar
postprocessing body at the stated visibility if and only if it is
a projective refinement of one orthonormal basis:
\[
 E_i=s_iP_{a(i)},\qquad s_i>0,\qquad
 \sum_{i:a(i)=a}s_i=1.
\]
Every other such POVM admits a pointwise increase in visibility.
\end{corollary}

\begin{proof}
The Haar body is full dimensional in the affine space \(\sum A_i=I\):
small bounded quadratic perturbations of constant responses realize
all small zero-sum Hermitian tuples by the Haar second moment.
A nontrivial supporting tuple has the normal form above.
The case \(\lambda=0\) forces a zero effect. When \(\lambda>0\),
the explicit optimal dual and complementary slackness force each
nonzero effect into one of its single basis-ray kernels.
Conversely \(W_i=P_{a(i)}\) supports the image of any projective
refinement, since both support and score equal \(H_d\).
For every remaining point, interiority and continuity permit a
small increase of \(\eta\).
\end{proof}

\subsection{The Euclidean selection inequality}
\label{geo:angles}

This subsection proves \eqref{geo:four-selection} and records all
constants used in the contact theorem. Put
\[
 f_t(\theta)=t\sin\theta-1-\cos\theta,\qquad
 g_t(\theta)=4t\sin\theta-2(1-t^2)(1-\cos\theta),\qquad
 H_t=g_t+4f_t.
\]

\paragraph{Planar balanced triples.}
Let three unit vectors in a plane enclose the origin, with cyclic gaps
\(\alpha,\beta,\gamma\in[0,\pi]\) summing to \(2\pi\).
For pair \(i,j\) and singleton \(k\), define \(w,v,A\) as in
\eqref{geo:pair-cell}, and set
\[
 D_k=|w|^2+t^2|v|^2+2t|w\times v|-A^2.
\]
All cyclic oriented sine terms are nonnegative. Expanding gives
\begin{equation}
 D_k=t^2(a_i-a_j)^2+a_ia_jg_t(\theta_{ij})
 +4a_k\{a_if_t(\theta_{ki})+a_jf_t(\theta_{jk})\}.
 \label{geo:planar-deficit}
\end{equation}
Take \(t=3/4\), \(\theta_0=2\arctan(1/t)\), and
\(\theta_1=2\arctan t=\pi-\theta_0\).
One has \(f_t\ge0\) on \([\theta_0,\pi]\) and
\(g_t\ge0\) on \([0,4\arctan t]\).
If two gaps are at least \(\theta_0\), choose their common vertex as
singleton; the remaining gap is at most \(4\arctan t\), so every
term in \eqref{geo:planar-deficit} is nonnegative.

Otherwise exactly one gap \(\alpha\) is at least \(\theta_0\), as
\(\theta_0<2\pi/3\). Its other gaps lie in
\([\theta_1,\theta_0]\). Add the deficits at the two endpoints \(i,j\)
of \(\alpha\):
\[
 D_i+D_j=t^2\{(a_j-a_k)^2+(a_i-a_k)^2\}
 +8a_ia_jf_t(\alpha)+a_ia_kH_t(\gamma)+a_ja_kH_t(\beta).
\]
On that middle interval,
\(\sin\theta\ge2t/(1+t^2)\) and
\(\cos\theta\le(1-t^2)/(1+t^2)\), hence
\[
 H_t(\theta)\ge
 \frac{4(t^4+3t^2-2)}{1+t^2}=\frac1{100}.
\]
Thus some deficit is nonnegative. Monotonicity in \(t\) gives
\eqref{geo:four-selection} for every balanced subset of at most
three rays, including limiting degeneracies.

\paragraph{A nearest-facet inequality.}
Now suppose four rays form a tetrahedron with the origin in its
interior; all other cases have a balanced subset of at most three.
Choose a facet of least plane distance \(h>0\) from the origin,
with outward normal \(e\), and write its vertices
\[
 n_i=he+Ru_i,\qquad R=\sqrt{1-h^2},\qquad u_i\in S^1\subset e^\perp.
\]
The radius-\(h\) ball lies in the tetrahedron and touches this facet
in its relative interior: equality with another facet would force
its unit normal to equal \(e\).
The projected triangle therefore encloses zero. Its cyclic gaps
\(\alpha,\beta,\gamma\) are in \((0,\pi)\), and
\[
 \Pi=\cos(\alpha/2)\cos(\beta/2)\cos(\gamma/2)>0,\qquad
 p=\frac{h^2}{R^2}.
\]
We claim
\begin{equation}
 \frac{h}{1-h}\le\sqrt{2\Pi},\qquad
 p\le\frac{2\Pi}{1+2\sqrt{2\Pi}}.
 \label{geo:facet}
\end{equation}
Write the fourth vertex as \(r-se\), \(0<s\le1\).
For an edge let \(\nu\) be its outward planar unit normal,
\(b=R\cos(\theta/2)>0\) its chord distance, and
\(\ell=2R\sin(\theta/2)\) its length. The adjacent facet equation is
\[
 (s+h)\nu\cdot x+(\nu\cdot r-b)z=sb+h\nu\cdot r.
\]
Its right side is positive, since the origin is in the interior.
Its distance from the origin is at least \(h\); squaring yields
\[
 \nu\cdot r\ge
 \frac{h(s+h)}{2b}-\frac{(s-h)b}{2h}.
\]
Multiply by \(\ell\), sum, and use the polygon identity
\(\sum\ell\nu=0\). This forces \(s>h\) and gives
\[
 h^2\frac{s+h}{s-h}\le\frac{\sum\ell b}{\sum\ell/b}
 =2R^2\Pi.
\]
The equality follows from the half-gap sum \(\pi\) and the identities
\(\sum\sin2A=4\sin A\sin B\sin C\),
\(\sum\tan A=\tan A\tan B\tan C\).
Since \(s\le1\), this gives \(h^2\le2(1-h)^2\Pi\), proving
\eqref{geo:facet}. In particular, concavity of \(\log\cos\) gives
\(\Pi\le1/8\) and \(h\le1/3\).

\paragraph{Lifting the planar deficit.}
Set \(t=9/10\). If \(w_0,v_0\) are the vectors for the projected
triangle, the normal component of \(w\times v\) has magnitude
\(R^2|w_0\times v_0|\), while
\[
 w=Rw_0+h(a_i+a_j-2a_k)e,\qquad
 v=Rv_0+h(a_i-a_j)e.
\]
Consequently
\[
 D_k\ge R^2D_k^{\rm planar}
 +h^2\{t^2(a_i-a_j)^2-8a_k(a_i+a_j)\}.
\]
Dropping the last positive square and using
\eqref{geo:planar-deficit},
\begin{equation}
 \frac{D_k}{R^2}\ge
 t^2(a_i-a_j)^2+a_ia_jg_t(\theta_{ij})
 +4a_k\{a_i(f_t(\theta_{ki})-2p)
             +a_j(f_t(\theta_{jk})-2p)\}.
 \label{geo:lifted}
\end{equation}
Call a gap large when it is at least $2\pi/3$. We prove the following
three bounds from \eqref{geo:facet}:
\begin{align}
 f_t(\theta)-2p&\ge0
 &&\text{for either large gap if two are large},
 \label{geo:coeff1}\\
 f_t(\alpha)-2p&\ge-1/200
 &&\text{for every }\alpha\ge2\pi/3,
 \label{geo:coeff2}\\
 H_t(\beta)-8p,\ H_t(\gamma)-8p&\ge1/25
 &&\text{if }\alpha\text{ is the unique large gap}.
 \label{geo:coeff3}
\end{align}

For \eqref{geo:coeff1}, put \(B=\theta/2\in[\pi/3,\pi/2]\),
\(z=\cos B\), and
\[
 r=\frac{\sqrt3}{2}\sin B-\frac z2,\qquad
 A=\frac{\sqrt3}{2}\sin B-z.
\]
The second large gap implies \(2\Pi\le zr\).
Both \(r,A\) increase with \(B\), while \(z\) decreases.
At \(B=\pi/3\), \(16rA^2=z\), so \(16rA^2\ge z\) throughout.
By \eqref{geo:facet} and monotonicity of \(u/(1+2\sqrt u)\),
\[
 p\le\frac{zr}{1+2\sqrt{zr}}\le zA
 =\frac12 f_{\sqrt3/2}(2B).
\]
The middle inequality is equivalent, by positive squaring, to the
inequality just proved. Since \(9/10>\sqrt3/2\),
\eqref{geo:coeff1} follows.

For \eqref{geo:coeff2}, put \(z=\cos(\alpha/2)\in[0,1/2]\).
The product of the other two half-gap cosines is at most
\((1-z)/2\). Thus
\[
 f_t(\alpha)-2p
 \ge2z\left[\frac9{10}\sqrt{1-z^2}-q(z)\right],\qquad
 q(z)=z+\frac{1-z}{1+2\sqrt{z(1-z)}}.
\]
On \(0\le z\le1/100\), use \(q(z)\le1\) and
\(\sqrt{1-z^2}\ge1-z^2\), obtaining
\[
 f_t(\alpha)-2p
 \ge-\frac{10009}{5000000}>-\frac1{200}.
\]
On \(1/100\le z\le1/4\), let \(k=\sqrt{z(1-z)}\).
The sign of \(q'\) is the sign of
\[
 N=(1-z)(4z-1)+4k^3,\qquad N'=5-8z+6k(1-2z)>0.
\]
Hence \(q\) has no interior maximum. Its endpoint values are at
most \(1/100+99/118<17/20\) and \(2/3<17/20\), whereas
\((9/10)\sqrt{1-z^2}\ge(9/10)\sqrt{15/16}>17/20\).
On \(1/4\le z\le1/2\), \(N\ge0\), so \(q(z)\le q(1/2)=3/4\);
meanwhile \((9/10)\sqrt{1-z^2}\ge9\sqrt3/20>3/4\).
These bounds prove \eqref{geo:coeff2}.

For \eqref{geo:coeff3}, the two middle gaps belong to
\([\pi/3,2\pi/3]\). For either such \(\theta\), maximizing the other
two half-gap cosine factors subject to one being large gives
\[
 p\le2\Pi\le\frac{\sqrt3\sin\theta-1-\cos\theta}{4}.
\]
Since
\(H_{9/10}=(36/5)\sin\theta-219/50-(181/50)\cos\theta\),
\[
 H_{9/10}(\theta)-8p\ge
 \left(\frac{36}{5}-2\sqrt3\right)\sin\theta
             -\frac{119}{50}-\frac{81}{50}\cos\theta.
\]
This expression is concave on the stated interval. Its smaller
endpoint value is \(18\sqrt3/5-619/100>1/25\), the last inequality
following after positive squaring from \(388800>388129\).

\paragraph{Completing selection.}
If two gaps are at least \(2\pi/3\), choose their common vertex
as singleton. The remaining gap is at most \(2\pi/3\), where
\(g_{9/10}\ge0\); \eqref{geo:coeff1} makes \eqref{geo:lifted}
nonnegative.
Otherwise let \(a,b\) be the amplitudes at the endpoints of the
unique large gap, and \(c\) the third amplitude. Add the endpoint
deficits, and use \eqref{geo:coeff2}--\eqref{geo:coeff3}:
\[
 \frac{D_i+D_j}{R^2}\ge
 \frac{81}{100}\{(a-c)^2+(b-c)^2\}
      -\frac{ab}{25}+\frac{c(a+b)}{25}
 =\frac{79}{200}(a+b-2c)^2
  +\frac{83}{200}(a-b)^2+\frac{c^2}{25}\ge0.
\]
Some deficit is nonnegative, proving \eqref{geo:four-selection}.
All amplitudes may vanish; degeneracies are covered either by
continuity or reduction to a smaller balanced subset.

\paragraph{The optional explicit cusp constant.}
For \(d=3\), the shifted diagonal Gram matrix on \(R_1\) is
\[
 M=\frac1{3888}
 \begin{pmatrix}251&-13&-13\\-13&20&-7\\-13&-7&20\end{pmatrix}.
\]
The exact Haar integrals give this matrix; off-diagonal phase
averages have larger unmultiplied coefficients, at least \(7/648\).
For its positive eigenvalues,
\(\lambda_{\min}(M)\ge4\det M/(\operatorname{tr}M)^2=325/338724\).
Also
\(\|z^*Dz-(7/36)\operatorname{tr}D\|_\infty\le(19/12)\|D\|\).
Dividing the \(L^2\) lower bound by this \(L^\infty\) upper bound
gives \(\kappa_3=325/536313\) in Lemma~\ref{geo:cusp}.

